% Trecho para inclusão no documento principal, após \begin{document}.
% Acrescentar ao preâmbulo: \usepackage{booktabs}
%
% No capítulo de implementação, inserir os seguintes labels imediatamente
% após os títulos correspondentes:
%   \subsection{Linux Network Namespaces}
%   \label{sec:linux_network_namespaces}
%
%   \subsubsection{Configuração da Topologia ICE para Hole Punching utilizando Linux Network Namespaces}
%   \label{sec:ice_topology_network_namespaces}
%
% Os demais labels referenciados já existem no capítulo fornecido.

\section{Ambiente experimental e cenários com PatchBay}
\label{sec:experimental_patchbay}

Os experimentos utilizam as topologias descritas nas
Seções~\ref{sec:basic_topology_network_namespaces}
e~\ref{sec:ice_topology_network_namespaces}, executadas com o
PatchBay sobre \textit{Linux Network Namespaces}. Esta seção
registra a configuração do ambiente e as diferenças entre os
cenários de teste. A organização das interfaces de transporte
e dos protocolos é apresentada no
Capítulo~\ref{chap:implementation}.

\subsection{Plataforma e serviços auxiliares}

Os testes foram executados em ambientes com Ubuntu 24 LTS
e Rocky Linux 8.10. A Tabela~\ref{tab:experimental_software}
apresenta as versões das dependências utilizadas.

\begin{table}[htbp]
    \centering
    \caption{Software utilizado no ambiente experimental.}
    \label{tab:experimental_software}
    \begin{tabular}{@{}p{0.22\textwidth}p{0.14\textwidth}p{\dimexpr0.64\textwidth-4\tabcolsep\relax}@{}}
        \toprule
        \textbf{Componente} & \textbf{Versão} & \textbf{Função} \\
        \midrule
        OpenSSL & 4.1.0 &
        Implementação QUIC e proteção criptográfica. \\
        \texttt{libnice} & 0.1.21 &
        Coleta de candidatos, negociação ICE e transporte
        direto ou via TURN. \\
        \texttt{coturn} & 4.6.3 &
        Fornecimento dos serviços STUN e TURN. \\
        PatchBay & 0.7.0 &
        Criação e configuração da rede de testes. \\
        \bottomrule
    \end{tabular}
\end{table}

Os serviços auxiliares executam no nó ISP, conforme a
topologia ICE apresentada na Seção~\ref{sec:ice_topology}.
Sua configuração é resumida na
Tabela~\ref{tab:experimental_services}.

\begin{table}[htbp]
    \centering
    \caption{Configuração dos serviços e da proteção QUIC.}
    \label{tab:experimental_services}
    \begin{tabular}{@{}p{0.28\textwidth}p{\dimexpr0.72\textwidth-2\tabcolsep\relax}@{}}
        \toprule
        \textbf{Elemento} & \textbf{Configuração} \\
        \midrule
        STUN/TURN &
        Comunicação por UDP, na porta 3478.
        O adaptador utiliza
        \texttt{NICE\_RELAY\_TYPE\_TURN\_UDP};
        TCP, TLS e DTLS estão desabilitados no servidor
        do laboratório. \\
        Autenticação TURN &
        Credenciais de longo prazo, com usuário e senha
        estáticos de laboratório e \textit{realm}
        \texttt{ice-test}. \\
        Sinalização &
        Servidor TCP na porta 5555, utilizado para trocar
        candidatos e credenciais ICE
        (\texttt{ice-ufrag} e \texttt{ice-pwd}). \\
        Proteção QUIC &
        TLS 1.3, com certificado autoassinado e chave
        privada no servidor. A correspondência entre
        certificado e chave é verificada na configuração. \\
        Verificação de identidade &
        O cliente utiliza \texttt{SSL\_VERIFY\_NONE},
        sem validar o certificado do servidor.
        Não há autenticação mútua por certificados. \\
        \bottomrule
    \end{tabular}
\end{table}

A configuração de certificados é restrita ao ambiente
experimental. Os testes exercitam o transporte protegido
pelo QUIC, mas não demonstram a validação da identidade
do servidor remoto.

\subsection{Topologia e cenários de comunicação}
\label{sec:experimental_scenarios}

A rede contém dois \textit{peers}, cada um atrás de um
roteador NAT, conectados por um roteador ISP. Os
\textit{peers} representam os \textit{gateways} da
arquitetura e executam tanto a implementação quanto as
aplicações TCP locais utilizadas nos testes de tunelamento.

Conforme as Figuras~\ref{fig:basic_topology}
e~\ref{fig:ice_topology}, a estrutura da rede é mantida
entre os cenários. As diferenças estão nas regras de NAT,
nos serviços auxiliares habilitados e na escolha do
transporte. As LANs são isoladas, embora utilizem o mesmo
prefixo \texttt{192.168.80.0/24}; a rede entre o ISP e os
roteadores utiliza \texttt{10.0.0.0/24}.

Os roteadores partem do perfil \texttt{Moderate} do
PatchBay. Nos cenários sem ICE, regras adicionais preservam
a porta de origem, permitindo o conhecimento prévio dos
endpoints públicos, como descrito na
Seção~\ref{sec:basic_topology_hole_punching}.
Nos cenários com ICE, a descoberta e a seleção dos
endpoints são realizadas pela \texttt{libnice},
conforme a Seção~\ref{sec:ice_library}.

A Tabela~\ref{tab:experimental_scenarios} organiza os
cenários por configuração de conectividade. Uma mesma
configuração pode ser exercitada por mais de um teste.

\begin{table}[htbp]
    \centering
    \small
    \caption{Cenários de comunicação avaliados com PatchBay.}
    \label{tab:experimental_scenarios}
    \begin{tabular}{@{}p{0.20\textwidth}p{0.38\textwidth}p{\dimexpr0.42\textwidth-4\tabcolsep\relax}@{}}
        \toprule
        \textbf{Cenário} &
        \textbf{Configuração} &
        \textbf{Objetivo} \\
        \midrule
        Comunicação sem ICE &
        UDP direto, endpoints previamente conhecidos
        e mapeamento de portas espelhado. &
        Exercitar Chownat e QUIC, com e sem tunelamento
        TCP, na topologia básica. \\
        Transporte ICE direto &
        STUN e sinalização habilitados, sem
        obrigatoriedade de relay. &
        Verificar a seleção de um caminho direto
        e a troca de datagramas pela interface
        \texttt{udp\_transport}. \\
        ICE com remapeamento de porta &
        Regras adicionais fazem a porta pública
        diferir da porta local. &
        Verificar a comunicação sem depender
        da preservação da porta de origem. \\
        Transporte ICE via TURN &
        TURN configurado e opção
        \texttt{force\_relay} habilitada. &
        Verificar o transporte de datagramas por
        um par de candidatos relay. \\
        Chownat sobre ICE &
        Tunelamento TCP com seleção ICE direta
        ou TURN obrigatório. &
        Verificar a integração do Chownat com
        ambos os caminhos de transporte. \\
        QUIC sobre ICE &
        Tunelamento TCP com seleção ICE direta
        ou TURN obrigatório. &
        Verificar a integração QUIC/ICE por meio
        do BIO descrito na
        Figura~\ref{fig:openssl_quic_bio_update}. \\
        \bottomrule
    \end{tabular}
\end{table}

Nos cenários diretos, os servidores auxiliam o
estabelecimento da comunicação, enquanto os dados da
aplicação trafegam entre os \textit{peers}. Nos cenários
TURN, os dados passam pelo servidor relay, conforme
a Figura~\ref{fig:turn_quic_tcp_tunneling}.

O uso de \texttt{force\_relay} torna obrigatório o caminho
TURN mesmo quando um caminho direto poderia existir.
Portanto, esses testes verificam a operação via relay,
sem representar, por si só, uma falha de conectividade
direta causada por um NAT mais restritivo.

Há também um teste prolongado de QUIC sobre TURN,
destinado a verificar a manutenção da sessão e a
renovação dos recursos TURN. A duração e a carga aplicada
são apresentadas na seção de cargas de trabalho e coleta
das medições.

\subsection{Parâmetros de execução}

A Tabela~\ref{tab:experimental_parameters} apresenta os
prazos e limites definidos pela implementação e pelas
configurações dos testes. Esses valores se aplicam às
etapas indicadas e não constituem um prazo global
de execução.

\begin{table}[htbp]
    \centering
    \caption{Parâmetros de execução dos testes de comunicação.}
    \label{tab:experimental_parameters}
    \begin{tabular}{@{}p{0.49\textwidth}p{0.10\textwidth}p{\dimexpr0.41\textwidth-4\tabcolsep\relax}@{}}
        \toprule
        \textbf{Parâmetro} &
        \textbf{Valor} &
        \textbf{Aplicação} \\
        \midrule
        Espera pela prontidão do ICE &
        5~s & Chownat/ICE e QUIC/ICE \\
        Leitura/escrita da sinalização &
        3~s & Chownat/ICE e QUIC/ICE \\
        Estabelecimento QUIC e abertura dos streams &
        5~s & Testes QUIC \\
        Espera por resposta de controle Chownat &
        1~s & Conexão e encerramento \\
        Limite de tentativas Chownat &
        5--20 & Com e sem ICE \\
        Inatividade para acionar \textit{keep-alive}
        da sessão &
        5~s & Camada \texttt{udp\_conn} \\
        \bottomrule
    \end{tabular}
\end{table}

No Chownat, o prazo de 1~s limita cada espera por uma
mensagem de controle; o limite de tentativas é cinco
nos testes com ICE e vinte nos testes sem ICE.
No QUIC, os 5~s abrangem o estabelecimento da conexão
e a abertura dos streams, compartilhando um único prazo
nessa etapa. Portanto, esses dois valores não representam
a mesma unidade de espera. São limites configurados na
implementação e nos testes, e não tempos de conexão medidos.

O intervalo de \textit{keep-alive} indicado corresponde
à camada de sessão descrita na
Seção~\ref{sec:api_and_event_loop}; os temporizadores
internos do ICE e do QUIC são gerenciados pelas respectivas
bibliotecas. O executor não configura explicitamente
atraso, perda de pacotes ou limitação de banda nos enlaces.

Como os \textit{namespaces} compartilham o mesmo
\textit{kernel} e os recursos da máquina hospedeira,
os resultados caracterizam esse ambiente local de
execução, conforme discutido na
Seção~\ref{sec:linux_network_namespaces}.

% Ajuste sugerido no capítulo de implementação: substituir a afirmação de
% que a topologia ICE herda "exatamente a mesma configuração" por:
%
% A topologia para ICE reutiliza a organização dos nós e
% o endereçamento descritos na
% Seção~\ref{sec:basic_topology_network_namespaces}.
% As regras adicionais de mapeamento espelhado dos testes
% sem ICE não são necessárias nessa configuração, pois
% os endpoints são descobertos e negociados pelo ICE.

\section{Resultados dos testes funcionais}
\label{sec:functional_results}

A avaliação funcional investigou três propriedades da integração:
a descoberta de um caminho entre gateways sem conhecimento prévio
das portas públicas, a preservação dos dados das aplicações TCP e
a manutenção e liberação dos recursos utilizados pelo transporte.
Para relacionar essas propriedades à implementação do
Capítulo~\ref{chap:implementation}, foram examinados os pacotes
capturados em diferentes interfaces, os pares selecionados pela
\texttt{libnice} e as verificações de conteúdo realizadas pelos testes.

Na coleta de 9 de novembro de 2026, em Rocky Linux 8.10, os 18 testes
da suíte regular e o teste prolongado QUIC/ICE/TURN foram aprovados.
Os resultados a seguir detalham as evidências dessa execução.

% ORIENTAÇÕES DE EDIÇÃO, não fazem parte da versão final do TCC:
% Substituir os seis placeholders pelos recortes indicados antes da entrega.
% Guia com arquivos, filtros, quadros e campos a expandir:
% docs/analysis/functional-pcap-figures-2026-11-09.md
% Evidências brutas: functional-results-2026-11-09/
% Os números de quadro só se aplicam aos arquivos originais dessa coleta.
% Não renumerar os quadros exportando apenas os pacotes filtrados.

\subsection{Estabelecimento da comunicação e caminhos selecionados}
\label{sec:functional_establishment}

A primeira diferença observável entre os transportes direto e ICE
é a obtenção do endpoint remoto. Nos testes sem ICE, Chownat e QUIC
utilizam os endereços previamente configurados e o mapeamento de
portas espelhado descrito na
Seção~\ref{sec:basic_topology_hole_punching}. A comunicação nesses
casos estabelece a referência funcional para a comparação com o
adaptador ICE. A descoberta e a negociação acrescentadas por esse
adaptador são examinadas pelo cenário com remapeamento de porta.

Nesse cenário, a captura da interface LAN do gateway A registra
um Binding Request originado em \texttt{192.168.80.20:54960}.
A mesma transação aparece na WAN do roteador A com origem
\texttt{10.0.0.2:38576}. A resposta STUN contém esse último endpoint
no atributo \texttt{XOR-MAPPED-ADDRESS}. A associação entre as duas
capturas é feita pelo identificador da transação STUN, preservado
durante a tradução. Portanto, a diferença entre as portas não é
apenas uma configuração declarada: ela aparece nos pacotes
efetivamente enviados e na resposta utilizada para a descoberta.

% FIGURA 1: teste network_ice_udp_transport_remap.
% Painel (a): pcap/peer_a/eth0.pcap, quadros 22 e 23.
% Painel (b): pcap/router-a/wan.pcap, quadros 22 e 23.
% Expandir UDP no Request e XOR-MAPPED-ADDRESS no Success Response.
% Destacar o mesmo Transaction ID nos dois painéis e as portas 54960/38576.
\begin{figure}[htbp]
    \centering
    \fbox{\parbox{0.90\linewidth}{\centering
        Inserir dois recortes da transação STUN: origem na LAN
        com porta 54960 e origem na WAN com porta 38576,
        incluindo o endereço informado na resposta.}}
    \caption{Descoberta do endpoint público do gateway A no cenário
    com remapeamento: correspondência entre a transação STUN na LAN
    e na WAN e o endereço retornado pelo servidor.}
    \label{fig:result_stun_remap}
\end{figure}

Após a troca de candidatos, a WAN do roteador A registra um Binding
Request destinado a \texttt{10.0.0.3:64027}, com
\texttt{USERNAME}, \texttt{MESSAGE-INTEGRITY} e
\texttt{USE-CANDIDATE}, seguido de uma resposta de sucesso.
O destino é o endpoint público do gateway B, cuja porta local
47643 foi traduzida para 64027. A seleção registrada pela
\texttt{libnice} corresponde ao par
\texttt{10.0.0.2:38576} e \texttt{10.0.0.3:64027}.
Essa sequência vincula a descoberta, a nomeação do par e a
comunicação posterior a portas públicas diferentes das portas locais.
Assim, nesse cenário, o transporte ICE remove a dependência do
mapeamento espelhado presente na referência sem ICE.

No teste QUIC/ICE direto, a captura permite acompanhar também a
passagem da negociação ICE para o estabelecimento QUIC. Os quadros
26 e 27 contêm a nomeação e a resposta STUN; posteriormente, o quadro
59 apresenta um datagrama UDP de 1200 bytes iniciado por um pacote
QUIC Initial. A sequência inclui Retry e pacotes Handshake entre
\texttt{10.0.0.2:49381} e \texttt{10.0.0.3:38368}.
Esses são os endpoints do par reflexivo selecionado nos logs.
Entre os checks e o Initial aparecem os pequenos datagramas de
controle do procedimento de descoberta mantido pelo backend QUIC.
A captura, portanto, permite distinguir a preparação do caminho
do handshake QUIC descrito na implementação.

% FIGURA 2: network_quic_ice_tcp/pcap/router-a/wan.pcap.
% Filtro: (frame.number >= 25 && frame.number <= 32) ||
%         (frame.number >= 59 && frame.number <= 73)
% Dois painéis: checks STUN (selecionar 26) e handshake (selecionar 59).
% Mostrar No., Time, Source, Destination, Protocol, Info e portas UDP.
% Sem chaves da sessão, não usar a imagem para afirmar validação de
% certificado nem para mostrar o conteúdo cifrado dos pacotes Handshake.
\begin{figure}[htbp]
    \centering
    \fbox{\parbox{0.90\linewidth}{\centering
        Inserir a sequência de checks STUN e, abaixo, os pacotes
        QUIC Initial, Retry e Handshake do mesmo par de endpoints.}}
    \caption{Estabelecimento QUIC sobre o caminho ICE direto:
    nomeação do par e pacotes QUIC observados entre os endpoints
    públicos selecionados.}
    \label{fig:result_ice_quic_direct}
\end{figure}

A Tabela~\ref{tab:functional_paths} sintetiza as seleções da suíte
regular. Nos quatro cenários ICE diretos, os dois candidatos foram
identificados como reflexivos de servidor. Nos três cenários com
TURN obrigatório, ambos os candidatos eram relay.

\begin{table}[htbp]
    \centering
    \small
    \caption{Caminhos observados nos testes de estabelecimento.}
    \label{tab:functional_paths}
    \begin{tabular}{@{}p{0.43\textwidth}p{\dimexpr0.57\textwidth-2\tabcolsep\relax}@{}}
        \toprule
        \textbf{Cenários} & \textbf{Caminho observado} \\
        \midrule
        Chownat e QUIC sem ICE & UDP direto, com endpoints configurados previamente. \\
        Transporte ICE direto e com remapeamento & Par reflexivo de servidor/reflexivo de servidor. \\
        Chownat/ICE e QUIC/ICE diretos & Par reflexivo de servidor/reflexivo de servidor. \\
        Transporte ICE, Chownat/ICE e QUIC/ICE via TURN & Par relay/relay, com allocations no servidor. \\
        \bottomrule
    \end{tabular}
\end{table}

O uso de TURN altera o caminho de dados, além da seleção informada
pela biblioteca. No teste QUIC/ICE/TURN, a resposta Allocate informa
\texttt{10.0.0.1:49191} como endereço de relay do gateway A.
O ChannelBind associa o canal \texttt{0x4000} ao relay remoto
\texttt{10.0.0.1:49172}; em seguida, a WAN registra mensagens
ChannelData entre \texttt{10.0.0.2:47158} e o servidor
\texttt{10.0.0.1:3478}. No quadro 64, uma dessas mensagens contém
1200 bytes de dados encapsulados, além do cabeçalho ChannelData de
quatro bytes. A captura evidencia, assim, o encapsulamento acrescentado
pelo adaptador ICE sem alterar a interface utilizada pelo backend QUIC
(Figura~\ref{fig:openssl_quic_bio_update}).

% FIGURA 3: network_quic_ice_turn_tcp/pcap/router-a/wan.pcap.
% Filtro: (frame.number >= 11 && frame.number <= 14) ||
%         frame.number == 30 || frame.number == 31 || frame.number == 64
% Destacar XOR-RELAYED-ADDRESS no 14; CHANNEL-NUMBER e XOR-PEER-ADDRESS
% no 30; cabeçalho ChannelData (canal 0x4000, Length 1200) no 64.
% No 64, o comprimento do payload UDP é 1204, sem contar o cabeçalho UDP.
\begin{figure}[htbp]
    \centering
    \fbox{\parbox{0.90\linewidth}{\centering
        Inserir os recortes de Allocate Success, ChannelBind e
        ChannelData, destacando endereços de relay e canal utilizado.}}
    \caption{Caminho QUIC/ICE via TURN: allocation, associação do canal
    ao relay remoto e transporte de dados por ChannelData.}
    \label{fig:result_turn_data_path}
\end{figure}

Esses cenários utilizam \texttt{force\_relay}. Demonstram o transporte
pelo servidor quando esse caminho é imposto, mas não avaliam a
transição automática de uma tentativa direta malsucedida para TURN.

\subsection{Transferência de dados e integração com aplicações TCP}
\label{sec:functional_data_transfer}

A verificação dos dados foi realizada nas interfaces entregues às
aplicações, pois a conclusão do handshake não demonstra que o túnel
preserva o conteúdo transmitido. Nos seis cenários de túnel TCP,
o gerador enviou 16384 bytes e comparou a resposta do serviço echo
byte a byte. O payload segue o padrão \(b_i = i \bmod 251\), incluindo
bytes nulos e valores fora do intervalo de caracteres imprimíveis.
Esse padrão permite detectar truncamentos decorrentes do tratamento
incorreto de dados binários como strings.

O cenário QUIC/ICE direto ilustra a relação entre as duas conexões
TCP locais e o transporte entre gateways. No gateway A, a aplicação
se conecta à porta 44444, e o quadro 86 da captura de loopback contém
os 16384 bytes entregues ao túnel. No gateway B, a implementação
abre outra conexão TCP com o serviço echo na porta 1234. Nessa
segunda conexão, o primeiro segmento de dados observado contém
1024 bytes, e os demais completam a mesma sequência de aplicação.
As fronteiras dos segmentos são diferentes, enquanto o conteúdo
entregue ao serviço é preservado.

% FIGURA 4: network_quic_ice_tcp, duas capturas de loopback.
% (a) pcap/peer_a/lo.pcap:
%     tcp.port == 56538 && tcp.port == 44444
%     selecionar quadro 86; mostrar TCP Len 16384 e bytes iniciais.
% (b) pcap/peer_b/lo.pcap:
%     tcp.port == 52972 && tcp.port == 1234
%     selecionar quadro 4; mostrar TCP Len 1024 e bytes iniciais.
% Em ambos, Follow TCP Stream em representação hexadecimal ajuda a
% mostrar o padrão 00 01 02 ...; a integridade total vem da comparação
% dos fluxos, e não apenas da coincidência desse prefixo.
\begin{figure}[htbp]
    \centering
    \fbox{\parbox{0.90\linewidth}{\centering
        Inserir dois painéis dos fluxos TCP locais: entrega do payload
        à porta 44444 no gateway A e recepção pelo echo na porta 1234
        no gateway B, com o início dos dados em hexadecimal.}}
    \caption{Integração com aplicações TCP no cenário QUIC/ICE direto:
    conexões locais distintas e preservação do conteúdo de aplicação,
    apesar das diferentes fronteiras dos segmentos.}
    \label{fig:result_tcp_application}
\end{figure}

Para conferir o conteúdo além do recorte visual, os dados foram
reconstruídos a partir dos números de sequência TCP dessas duas
capturas. Os quatro fluxos unidirecionais, correspondentes à ida e
à volta em cada gateway, continham exatamente 16384 bytes,
sem lacunas, e eram idênticos ao payload de referência.
Isso relaciona os bytes enviados pela aplicação aos recebidos
pelo serviço remoto e aos devolvidos na resposta. A evidência
sustenta o tunelamento de dados descrito na
Figura~\ref{fig:udp_conn_tcp_tunneling_diagram}: os segmentos TCP
pertencem às conexões locais, e os dados atravessam a rede externa
pelo transporte QUIC selecionado.

% Reconstrução e SHA-256 dos quatro fluxos:
% functional-results-2026-11-09/tcp-reassembly.json
% Todos: 4348e3b98e8a327b34ced39c1da9e67cdb4cd5e48e4d7960607a3ae403d35f0c
% O tamanho 16384 no loopback não é uma medida da MTU do enlace WAN.

\begin{table}[htbp]
    \centering
    \small
    \caption{Resultados das verificações de conteúdo e entrega às aplicações.}
    \label{tab:functional_transfer}
    \begin{tabular}{@{}p{0.36\textwidth}p{\dimexpr0.64\textwidth-2\tabcolsep\relax}@{}}
        \toprule
        \textbf{Grupo de testes} & \textbf{Resultado observado} \\
        \midrule
        Chownat e QUIC sem túnel & Três rodadas de requisição/resposta com conteúdo e sequência conferidos por protocolo. \\
        Transporte ICE direto, remapeado e via TURN & Mensagens recebidas nos dois sentidos com conteúdo e tamanho esperados. \\
        Chownat com túnel TCP & Echo de 16384 bytes íntegro nas variantes sem ICE, ICE direto e ICE/TURN. \\
        QUIC com túnel TCP & Echo de 16384 bytes íntegro nas variantes sem ICE, ICE direto e ICE/TURN. \\
        \bottomrule
    \end{tabular}
\end{table}

Nas quatro integrações de túnel com ICE, os contadores dos peers
também registraram 16384 bytes enviados e recebidos por extremo e
um único evento de conexão estabelecida. O resultado simultâneo
nas variantes direta e TURN indica que a escolha do caminho ficou
encapsulada no transporte, mantendo a entrega esperada pelas
aplicações TCP locais. Essa conclusão se refere às cargas e às
condições de rede exercitadas nesta coleta.

\subsection{Manutenção e encerramento das conexões}
\label{sec:functional_maintenance_shutdown}

Uma execução curta pode concluir a transferência antes de ser
necessário renovar os recursos TURN. Por isso, o cenário prolongado
manteve a sessão QUIC/ICE/TURN durante 67 transferências, com
intervalos de 10~s entre os blocos a partir do segundo envio.
Cada peer confirmou 1097728 bytes enviados e recebidos, e o teste
completo, incluindo preparação e encerramento, terminou em 663,81~s.
Sua duração permitiu observar renovações durante a mesma sessão.

A captura do gateway A mostra novos pedidos CreatePermission
aproximadamente aos 241 e 481~s, ambos respondidos com sucesso.
Aos 540~s, aparecem um Refresh e um ChannelBind, também seguidos
de respostas de sucesso; a resposta ao Refresh informa
\texttt{LIFETIME=600}. A
Tabela~\ref{tab:functional_lifecycle} situa essas transações na
captura, e a Figura~\ref{fig:result_turn_renewals} reserva os
recortes correspondentes. Os registros do servidor confirmam
as mesmas operações para a sessão do outro gateway.

\begin{table}[htbp]
    \centering
    \small
    \caption{Transações de manutenção e remoção na captura do gateway A
    durante o teste prolongado. Tempos relativos ao início dessa captura.}
    \label{tab:functional_lifecycle}
    \begin{tabular}{@{}p{0.15\textwidth}p{0.24\textwidth}p{\dimexpr0.61\textwidth-4\tabcolsep\relax}@{}}
        \toprule
        \textbf{Tempo (s)} & \textbf{Quadros} & \textbf{Transação e resposta} \\
        \midrule
        241,29 & 1312--1313 & CreatePermission e resposta de sucesso. \\
        481,73 & 2456--2457 & Nova renovação de permissão e resposta de sucesso. \\
        540,54 & 2757--2760 & Refresh com resposta LIFETIME 600 e renovação ChannelBind. \\
        662,14 & 3371--3374 & Remoção com LIFETIME zero, erro 438, novo pedido e sucesso. \\
        \bottomrule
    \end{tabular}
\end{table}

% FIGURA 5: teste prolongado, pcap/peer_a/eth0.pcap.
% Filtro: (frame.number >= 1312 && frame.number <= 1313) ||
%         (frame.number >= 2757 && frame.number <= 2760)
% Mostrar Time, tipo STUN e correspondência dos Transaction IDs.
% Expandir LIFETIME=600 no quadro 2759. Não confundir esse valor
% de validade da allocation com a duração medida do teste.
\begin{figure}[htbp]
    \centering
    \fbox{\parbox{0.90\linewidth}{\centering
        Inserir as transações CreatePermission, Refresh e ChannelBind
        da sessão prolongada, incluindo as respostas e o LIFETIME
        da allocation renovada.}}
    \caption{Renovação dos recursos TURN durante a sessão QUIC/ICE:
    permissões, allocation e canal mantidos antes do encerramento.}
    \label{fig:result_turn_renewals}
\end{figure}

No encerramento, a captura revela uma condição adicional ao
fechamento normal. O gateway A envia Refresh com
\texttt{LIFETIME=0}, mas recebe erro 438 (\textit{Stale Nonce}).
Em seguida, envia outro Refresh com o nonce fornecido nessa
resposta e recebe sucesso com tempo de vida zero.
Os dois pares de mensagens possuem identificadores de transação
distintos: trata-se de um novo pedido autenticado, e não apenas
da retransmissão do datagrama anterior.

% FIGURA 6: teste prolongado, pcap/peer_a/eth0.pcap.
% Filtro: frame.number >= 3371 && frame.number <= 3374
% Expandir: LIFETIME=0 no 3371; ERROR-CODE=438 e NONCE no 3372;
% NONCE e LIFETIME=0 no 3373; LIFETIME=0 no Success Response 3374.
% O NONCE do 3373 é igual ao do erro 3372 e diferente do pedido 3371.
% Transaction IDs: 3371/3372 = 4127ab6b9d69af4eeb2ad43d;
%                  3373/3374 = c564076ffa9e24560e8400c6.
\begin{figure}[htbp]
    \centering
    \fbox{\parbox{0.90\linewidth}{\centering
        Inserir os quatro pacotes do fechamento TURN:
        Refresh com tempo de vida zero, erro 438,
        novo Refresh com nonce atualizado e resposta de sucesso.}}
    \caption{Remoção da allocation TURN após expiração do nonce:
    tratamento do erro 438 e conclusão da nova transação de remoção.}
    \label{fig:result_turn_stale_nonce}
\end{figure}

Essa sequência permite verificar o efeito da correção aplicada
à \texttt{libnice}: o recebimento do erro não encerrou o procedimento
antes da remoção efetiva. O servidor registrou o erro 438 e a
remoção com tempo de vida zero para os dois peers, e o verificador
encerrou com os registros:

\begin{quote}
\begin{verbatim}
TURN_RENEWED allocations=2 channels=2
TURN_RELEASED allocations=2
\end{verbatim}
\end{quote}

O fechamento do túnel TCP é observado separadamente nas capturas
locais da Figura~\ref{fig:result_tcp_application}: os quadros
120 e 122 do gateway A e 39 e 40 do gateway B contêm FIN.
O serviço echo chegou a EOF, os peers terminaram e as fixtures
de integração confirmaram o estado fechado. No teste Chownat sem
túnel dedicado ao encerramento, foram conferidas 20 mensagens
recebidas no cliente e 21 no servidor antes da liberação da sessão.

O teste prolongado demonstra a continuidade da transferência durante
as renovações observadas e a conclusão do encerramento após erro 438.
Ele não separa a contribuição de cada mecanismo de manutenção:
o keep-alive da sessão, os temporizadores QUIC e as mensagens STUN
permaneceram habilitados durante a execução.

\subsection{Comportamento diante de falhas}
\label{sec:functional_failures}

Os testes locais de falha examinam o caminho de erro do adaptador ICE.
Para considerar cada caso aprovado, foi exigido que o transporte
retornasse \texttt{UDP\_TRANSPORT\_ERROR}, atingisse o estado
\texttt{FAILED} e liberasse seus recursos. O simples término do
processo pelo limite de tempo do executor não atende a esse critério.

\begin{table}[htbp]
    \centering
    \small
    \caption{Condição introduzida e evidência do tratamento de falha.}
    \label{tab:functional_failures}
    \begin{tabular}{@{}p{0.32\textwidth}p{\dimexpr0.68\textwidth-2\tabcolsep\relax}@{}}
        \toprule
        \textbf{Condição} & \textbf{Evidência observada e interpretação} \\
        \midrule
        Sinalização indisponível & Conexão recusada na callback de sinalização; falha sem atingir READY, antes da troca de candidatos. \\
        Credenciais TURN inválidas & Respostas 401, nenhuma allocation criada e nenhum candidato relay; falha com relay obrigatório. \\
        Peer ausente & Candidato remoto configurado sem respondedor; verificações terminam em FAILED, sem atingir READY. \\
        Nonce expirado no fechamento & Allocation criada com credenciais válidas; erro 438 seguido de novo pedido e remoção confirmada no servidor. \\
        \bottomrule
    \end{tabular}
\end{table}

% Evidência de falhas: short/LastTest.log, procurar os nomes dos testes
% ice_failure_signaling_unavailable, ice_failure_turn_bad_credentials,
% ice_failure_remote_absent e ice_failure_turn_stale_nonce.
% Esses casos locais não foram capturados pelo runner PatchBay.
% Não usar uma imagem de um caso de rede bem-sucedido como evidência
% de credenciais inválidas. Um primeiro 401 também ocorre no desafio
% normal de autenticação, conforme quadros 11--14 da FIGURA 3.

O erro 401, isoladamente, não distingue credenciais inválidas do
desafio inicial da autenticação TURN. Na execução bem-sucedida da
Figura~\ref{fig:result_turn_data_path}, o primeiro Allocate também
recebe 401, mas o pedido seguinte obtém sucesso. No teste negativo,
a evidência relevante é a ausência de qualquer allocation bem-sucedida,
combinada à ausência de candidato relay e ao estado FAILED.
Essa distinção evita interpretar uma etapa normal de autenticação
como falha do transporte.

Nos quatro casos locais, a contagem de descritores de arquivo após
o fechamento retornou ao valor anterior à abertura do adaptador.
Esse critério é verificado pelas fixtures, após a inicialização
dos recursos globais da GLib, e aparece no registro de aprovação
como \texttt{expected failure and cleanup confirmed}.
Os três cenários de indisponibilidade, credenciais inválidas e
peer ausente terminaram antes do limite de 8~s da fixture, que
utiliza temporizadores ICE reduzidos exclusivamente nesses testes.

O caso de nonce expirado complementa a sessão prolongada:
reproduz localmente a expiração antes do fechamento e exige
\texttt{server removal confirmed}. A captura da
Figura~\ref{fig:result_turn_stale_nonce} fornece a comprovação
em pacotes do mesmo tratamento no cenário de rede prolongado.
O alcance desses resultados é o estabelecimento e a liberação
do transporte; interrupções de sessões já estabelecidas e
recuperação após mudança de rede permanecem fora desta avaliação.
